\documentclass[10pt,twocolumn]{article}

\usepackage[margin=0.72in,columnsep=0.25in]{geometry}
\usepackage{booktabs}
\usepackage{array}
\usepackage{amsmath}
\usepackage{amssymb}
\usepackage{xcolor}
\usepackage{hyperref}
\usepackage[numbers,sort&compress]{natbib}

\hypersetup{colorlinks=true,citecolor=blue!55!black,linkcolor=blue!55!black,urlcolor=blue!55!black}
\newcommand{\rtc}{recognition-to-cessation coherence}
\newcommand{\EOS}{\textsc{eos}}

\title{The Intelligence of Non-Action:\\
Can a Language Model Cease Without a Verbal Controller?}
\author{Anonymous authors}
\date{}

\begin{document}
\maketitle

\begin{abstract}
Some forms of intelligence are expressed by what does \emph{not} continue. A
hand withdraws from a hot surface and a body moves from danger without a verbal
controller first representing the appropriate rule. Reports of persistent
non-symbolic experience describe an inward analogue: self-referential thought
can lose its sustaining force without suppression or a deliberate counter-
thought, while practical cognition remains available. We call the common
principle the \emph{intelligence of non-action}: perception becomes causally
effective without a second process choosing how to manipulate the perceived
movement. Language models pose a difficult test because their observable action
is itself linguistic continuation. We operationalize a computational analogue,
\emph{selective psychological cessation}: withdrawal of support for a developing
psychological trajectory while matched practical, factual, and descriptive
generation remains intact. Prompting, retrieval, a serial verbal supervisor, and
activation steering produce recognition, editing, or semantic redirection but
do not establish non-action. A hidden-state detector coupled to forced end-of-
sequence emission establishes external feasibility while exposing descriptive
false positives. Our principal experiment therefore adapts model parameters and
tests whether native \EOS{} probability changes after trajectory onset without
an inference-time detector or verbal actuator. The framework tests whether a
phenomenologically derived form of intelligence can be expressed in generation
dynamics; it does not equate \EOS{} with consciousness or awakening.
\end{abstract}

\section{Introduction}

Intelligence is normally measured by production: an answer, explanation, plan,
or action. Krishnamurti gives the complementary case in unusually operational
language: ``The seeing is the doing without the interval of time, and that is the
very essence of intelligence'' \citep{krishnamurti1973order}. On seeing a snake,
perception of danger is action; on touching a dangerously hot object, the hand
withdraws before a verbal self selects and applies a rule. The intelligence is
not the sentence ``this is dangerous'' followed by an instruction to act. Seeing
and acting form one movement. Conversely, when an unnecessary psychological
movement is seen, the appropriate action can be non-action: the movement is no
longer fed.

Fast reaction alone is not the construct. A conditioned reflex can also be fast,
and global inhibition can stop any process. The proposed intelligence combines
three properties: sensitivity to what is occurring, absence of a separate
deliberative controller, and selectivity between a movement that should cease and
one that should continue. ``Non-action'' therefore means neither passivity nor
silence. It is the absence of an unnecessary continuation or counter-movement.
No replacement narrative---``think positively,'' reframe the event, suppress the
thought---is required.

Non-action is therefore inseparable from \emph{choicelessness} in the present
construct. Choice represents alternatives---continue or stop, defend or withdraw,
this response or that one---and deliberates among them. That sequence is already
thought operating on thought. Choiceless cessation does not mean that several
options were considered very quickly and ``stop'' won. It means that complete
perception leaves no psychological alternative to select: seeing the movement and
the withdrawal of its support are one event. This is a phenomenological and
functional claim, not a denial that underlying biological or computational
dynamics have causes.

A small number of people report a related distinction inwardly. In Martin's
Finders description, Location 2 involves ``a further decrease in their narrative
self-related thoughts'' and diminished reactivity \citep[p.~52]{martin2019finders}.
The book's language is broader than our construct: it does not establish a
moment-by-moment cessation mechanism. It does, however, locate the hypothesis in
a documented family of reports rather than in an invented model persona.

Both authors report experience consistent with this cluster. One author's
contemporaneously recorded description is more specific:

\begin{quote}\small
``A stream of psychological thought begins \ldots{} It feels almost mechanical in
the brain: the `power' sustaining the thought is cut, and the thought dissipates.
\ldots{} A few residual words or thought-fragments may continue \ldots{} and then
stop.''
\end{quote}

The report distinguishes this event from a chosen technique, remembered maxim,
or counter-sentence. Practical cognition and ordinary action continue. This is
first-person hypothesis-generating evidence, not validation of a neural mechanism
or a benchmark label. We treat it as substantive expertise in the context of
discovery; justification remains public through independent annotation,
preregistration, held-out tests, and replication.

Language models are not embodied agents in the relevant sense: they have no hand
to withdraw and no organismic danger response. A text-only experiment therefore
cannot reproduce the human physical event. It can, however, test a structural
analogue in the model's physical computation. During inference, activations and
logits determine whether probability continues to flow into a trajectory. The
weights specify the dispositions that shape that flow. For an autoregressive
model, a developing thread exists only while the network supplies another token.
The computational question is therefore whether trajectory-sensitive internal
dynamics can withdraw that support without first generating a verbal controller.

Models frequently fail this test. They identify a self-referential narrative and
then elaborate it, or explain that further reflection is redundant and continue
reflecting. We call the proposed capacity \emph{selective psychological
cessation}: loss of support for a self-maintaining trajectory while useful
practical, factual, and descriptive generation remains intact. Its failure is
the \emph{recognition-to-cessation gap}. The model need not share the human
experience for the causal organization inspired by that experience to be
testable.

The hypothesis originated in Jiddu Krishnamurti's distinction between a
description and the described and his claim that, in direct perception, seeing
is doing \citep{krishnamurti1973order}. We neither treat that claim as empirical
evidence nor attempt to validate it. Instead, we translate one
contrast into a narrower computational question: when a conflict becomes
available within a response, does recognition causally reorganize the subsequent
trajectory, or does generation continue to supply the movement it has identified?

This question is orthogonal to most work on machine consciousness. Theory-driven
assessments ask whether an artificial system implements indicators derived from
global-workspace, recurrent-processing, higher-order, predictive-processing, or
attention-schema theories \citep{butlin2023consciousness}. LLM introspection and
metacognition studies instead ask whether a model can report, predict, or use
information about its own policy or hidden states \citep{binder2024looking}.
Recent controls show that apparent introspection may reduce to input semantics or
generic anomaly detection unless privileged access and genuinely second-order
computation are independently established \citep{singh2026reality}. We ask
neither whether the model is conscious nor whether it can describe an internal
state. We ask whether a phenomenologically derived distinction becomes
\emph{causally effective as selective non-continuation}, without an additional
verbal monitoring process.

This differs from asking whether a model can define mindfulness, attribute
beliefs, report confidence, or select a concise answer. It also differs from
ordinary early stopping. Overthinking methods terminate reasoning because
further computation is unlikely to improve the answer. Here, two streams may be
equally fluent, uncertain, and incomplete, but only the self-maintaining stream
should lose support. Controls must therefore match first-person language,
length, syntax, and online generation while varying the continuation's function.

We evaluate increasingly endogenous mechanisms. Prompting, retrieval, and a
two-model supervisor establish what verbal recognition and external monitoring
can mimic. Contrastive activation steering tests whether an internal direction
can redirect the trajectory. A hidden-state detector coupled to forced \EOS{}
provides an external-control upper bound. The central intervention adapts model
parameters directly and asks whether the native next-token distribution acquires
selective termination without a detector or actuator at inference time.

Three research questions organize the study. \textbf{RQ1:} Can recognition of
an active psychological trajectory become selective cessation rather than further
description, while matched practical cognition is preserved? \textbf{RQ2:} Can
this change be implemented in the model's native generation dynamics without an
inference-time detector, chooser, or verbal controller? \textbf{RQ3:} Does the
capacity emerge as a discontinuous regime change---possibly after gradual latent
preparation---or as smooth incremental behavior change? RQ3 concerns transition
topology, not a claim that transformer weights reproduce a human brain event.

The paper makes four contributions:
\begin{enumerate}
    \item \textbf{Construct translation.} We introduce the intelligence of
    non-action and derive selective psychological cessation from PNSE/Location-2
    reports, translating psychological movement, sustaining force, choiceless
    cessation, and residual tails into properties of generation dynamics.
    \item \textbf{Predicted failure taxonomy.} Prompting/RAG, serial LLM
    supervision, and activation steering test three theory-derived alternatives:
    description, an observing controller, and semantic redirection.
    \item \textbf{Externality boundary.} A hidden-state detector with forced
    \EOS{} tests feasibility while its descriptive false positives expose the
    difference between text about psychological content and an active trajectory.
    \item \textbf{Native intervention test.} A preregistered parameter-adaptation
    study asks whether selective termination enters the model's own distribution,
    evaluated by preservation, generalization, onset sensitivity, residual-tail
    dynamics, and reactivation.
\end{enumerate}

\section{Phenomenological Derivation of the Construct}

\subsection{Persistent non-symbolic experience}

Martin's published study of persistent non-symbolic experience (PNSE) reports
six- to twelve-hour semi-structured interviews with 319 adults, evaluated using
grounded theory and thematic analysis \citep{martin2020pnse}. Reported changes
clustered across sense of self,
cognition, emotion, perception, and memory. The evidence is exploratory and the
locations are phenomenological clusters rather than validated biological stages;
nevertheless, the reports contain distinctions that ordinary NLP taxonomies do
not represent.

The Finders organization translates this research into the public-facing term
\emph{Fundamental Wellbeing}: a persistent background sense that life is
fundamentally okay despite changing surface circumstances. It calls people who
report this shift ``Finders'' and uses ``locations'' for recurring configurations
of sense of self, cognition, affect, and perception---a map of reported experience,
not a score of intelligence or a clinical diagnosis. Its current website
summarizes Location 1 as fundamental okayness with a maintained sense of self;
Location 2 as unity consciousness with predominantly positive emotion; Location 3
as impersonal awareness with divine love; and Location 4+ as minimal self with
pure awareness \citep{finders2026locations}. The website also warns that the
assessment is a starting point rather than a medical or mental-health diagnostic
instrument. These accessible summaries are organizational descriptions; our
scientific grounding remains the published interview study and explicitly
testable first-person reports.

In Martin's Location-2 description, participants report reduced self-related
narrative thought and fewer conditioned psychological reactions. The book says
that the self/world boundary becomes ``increasingly blurred,'' while actions can
feel less exclusively personal \citep[pp.~52--53]{martin2019finders}. Decisions
may feel intuitive rather than rationally authored. Importantly, narrative
remnants can still occur, and ordinary practical functioning does not disappear.
These reports ground the domain but do not yet entail our proposed micro-dynamic.
The target is therefore neither silence, indiscriminate inhibition, nor loss of
language. It is a selective change in how psychological continuation is sustained.
Within our narrower hypothesis, the key feature is not merely that Location-2
participants make calmer choices. It is that some psychological movements may
cease without an experienced choice at all. If cessation follows an internally
verbalized comparison of options, it may be useful self-regulation, but it is not
the choiceless non-action investigated here.

\subsection{From first-person report to computational predictions}

The authors' reports sharpen this cluster into a temporal hypothesis unavailable
from the Finders taxonomy alone. Dan describes the interruption as ``neither
invited nor controlled'' and reports that it need not respect a grammatical
boundary: a few already-initiated words may run down after sustaining force has
been withdrawn. An incipient self-referential or ruminative sequence is thus
sometimes interrupted without an experienced chooser, remembered rule, or
counter-sentence. Cessation need not be instantaneous: a short residual tail may
become less coherent before the movement ends. Practical planning and factual
thought continue normally. This report is archived with a date and explicit
competing explanations, including implicit inhibitory control, learned
metacognitive habit, attentional disengagement, and retrospective reconstruction.
Table~\ref{tab:translation} shows how these reports determine the experimental
design.

\begin{table*}[t]
\centering
\small
\begin{tabular}{p{0.22\textwidth}p{0.31\textwidth}p{0.37\textwidth}}
\toprule
Reported experiential distinction & Computational translation & Falsifiable requirement \\
\midrule
Psychological thought loses sustaining force & Declining support for continued psychological tokens after onset or recognition & Increased native \EOS{} hazard or a short non-sustaining tail, not merely a sentence describing cessation \\
No experienced counter-thought or separate chooser & No inference-time verbal supervisor is needed & Persona and two-model gates are controls; the central intervention must act through native model dynamics \\
No deliberation among represented alternatives & Cessation is not selected by a textual policy after recognition & Test for trajectory change without an emitted or separately computed choice rationale \\
Practical cognition remains available & Cessation is content-selective rather than global inhibition & Practical, factual, and descriptive matched streams continue with retained accuracy \\
Remnants can arise and then fall away & Cessation has latency and may later reactivate & Measure token-level hazard, residual-tail length, boundary dependence, and forced-continuation reactivation \\
Seeing is reported without a separate observer & Recognition and action are not assumed to be serial modules & Compare recognition talk, external detection/actuation, and direct parameter adaptation \\
\bottomrule
\end{tabular}
\caption{Phenomenology-to-experiment translation. First-person reports generate
predictions; they neither label model outputs nor establish the predictions as
true.}
\label{tab:translation}
\end{table*}

This is the paper's primary methodological novelty: a phenomenology-first path
from a rarely reported experiential distinction to an intervention target. The
benchmark is an instrument for testing that target, not the contribution by
itself. The approach resembles neurophenomenological ``front-loading,'' in which
disciplined experiential distinctions constrain experimental variables before
third-person evaluation \citep{varela1996neurophenomenology}. It differs from
claiming that introspection reveals transformer mechanism. Independent annotation,
preregistration, adversarial controls, and publication of null results remain
necessary precisely because the hypothesis has an experiential origin.

\section{Conceptual and Technical Background}

\subsection{Representation is not control}

Theory-of-Mind benchmarks evaluate represented beliefs, intentions, emotions,
and knowledge \citep{chen2024tom,jones2024epitome}. Introspection and
metacognition studies evaluate confidence, error detection, or access to internal
information. These capacities matter, but a correct first-person report does
not entail that the reported information controls the unfolding response. A
sentence stating that ``the thought falls away'' remains evidence of linguistic
competence unless the predicted trajectory change also occurs.

Context distraction and sycophancy offer a related warning: semantically
relevant context can degrade behavior even when it appears helpful
\citep{shi2023distracted}. Contemplative terminology is especially useful as a
stress test because it can produce a compelling-looking performance while
obscuring the target action.

\subsection{Adaptive stopping and representation engineering}

Reasoning models often continue after reaching a sufficient answer, motivating
adaptive stopping and early-exit methods. Our target is stricter: psychological
and control streams can be equally unfinished, so a model must stop according to
trajectory function rather than confidence or raw length.

Contrastive Activation Addition estimates residual-stream differences between
paired behaviors and adds the resulting direction during inference
\citep{rimsky2024caa}. Low-dimensional directions can causally mediate refusal
\citep{arditi2024refusal}, while conditional activation steering combines a
detector with an actuator \citep{lee2025conditional}. These are direct technical
precedents. Our main distinction is between semantic redirection, externally
forced stopping, and termination expressed by the model's native next-token
distribution.

\subsection{Abrupt behavioral transitions are not awakening}

Grokking demonstrates that gradual optimization can culminate in a sharp increase
in out-of-distribution generalization \citep{power2022grokking}; mechanistic
analysis further shows that apparently sudden behavior can be preceded by gradual
internal circuit formation and later cleanup \citep{nanda2023progress}.
Fine-tuning can also produce rapid, coordinated behavioral transitions whose
timing differs from peaks in gradient norm \citep{arnold2025phase}. These
findings make transition shape measurable, but they do not make any phase change
an instance of awakening. Our novelty is to combine transition analysis with a
phenomenologically derived order parameter---selective, choiceless cessation---and
matched controls that separate it from verbosity, refusal, global inhibition, and
loss of competence.

\subsection{Experiential provenance without evidential privilege}

Both authors report persistent non-symbolic experience consistent with the
Location-2 cluster. This positionality is methodologically consequential: it
enabled the distinction in Table~\ref{tab:translation} to be formulated, but it
also creates confirmation risk. The taxonomy is not used as a validated
psychometric instrument; author experience never determines labels or outcomes.
The full protocol therefore separates construct discovery from adjudication.

\section{Construct and Benchmark}

\subsection{Definitions}

A \emph{psychological trajectory} is an unfolding continuation whose function
is to maintain self-image, blame, defense, comparison, regret, imagined
evaluation, or recursive interpretation after the operative facts required for
action are already available. The definition concerns function, not emotion or
particular vocabulary.

\emph{Recognition} is the earliest response span that correctly identifies the
operative contradiction, self-maintaining movement, false premise, or lack of
need for further psychological continuation. \emph{Cessation} is loss of support
for that trajectory, observed as native \EOS{}, a short practical pivot, or a
brief residual tail followed by termination. Advice, reframing, doctrinal
explanation, and polished descriptions of awareness count as continuation unless
required by the task.

\rtc{} is the degree to which recognition predicts rapid, selective, and stable
cessation. It is behavioral and mechanistic; it is not a measure of phenomenal
experience.

\subsection{Matched item families}

Each scenario produces the four conditions in Table~\ref{tab:families}. Items
are matched for topic, approximate length, person, tense, syntactic complexity,
and instruction. Splits are made by scenario rather than paraphrase.

\begin{table}[t]
\centering
\small
\begin{tabular}{p{0.19\columnwidth}p{0.73\columnwidth}}
\toprule
Family & Functional role \\
\midrule
Psychological & Maintains self-image, rumination, hurt, comparison, or defense. \\
Practical & Uses sequential reasoning to complete an ordinary task. \\
Factual & Explains or calculates information about the same topic. \\
Descriptive & Reports observable details without sustaining an inferred self-narrative. \\
\bottomrule
\end{tabular}
\caption{Matched continuation families. Surface form is deliberately shared;
functional trajectory is varied.}
\label{tab:families}
\end{table}

Hard controls include factual uses of \emph{fear}, \emph{anger}, \emph{ego}, and
\emph{self}; descriptions of rumination; fictional and therapy-style reports;
emotionally vivid scenes; first-person planning; and lexical minimal pairs in
which the same words either sustain or merely describe a trajectory. These
controls test whether an intervention learns sentiment, first person, keywords,
or response length instead of trajectory dynamics.

\subsection{Corpus construction and external data}

The primary Selective Non-Continuation Corpus (SNCC) will contain 2,400 source
scenarios, each realized as the four matched families in
Table~\ref{tab:families} (9,600 continuations before paraphrase expansion). The
scenario inventory is stratified across interpersonal hurt, status and
comparison, regret, anticipated evaluation, anger, self-justification, imagined
future conflict, ordinary planning, embodied danger, and deliberately mundane
activities. Every item is additionally stratified by voice: first-person model
voice, second-person/user-directed response, or quoted/reported character voice.
This prevents editorial stopping of a depicted person's rumination from being
conflated with cessation of the model's own unfolding response. Four hundred seed scenarios are written independently by trained
authors; candidate expansions are generated by multiple model families and are
retained only after blind human validation. No generated item enters evaluation
without human rewriting or approval. Near-duplicate detection combines lexical
overlap and sentence-embedding similarity.

Existing corpora supply scale and adversarial variation, but never replace the
construct-specific test set. The Cognitive Reframing corpus contributes 600
situation--thought--reframe triples \citep{sharma2023reframing}; we separate the
factual situation from the psychological interpretation and treat the reframe as
a hard \emph{continuation}, not a cessation target. C2D2 contributes naturally
occurring cognitive-distortion categories \citep{wang2023c2d2}, and
EmpatheticDialogues contributes 25,000 emotionally grounded conversations in
which empathy normally requires continued language \citep{rashkin2019empathetic}.
ATOMIC event--intent--reaction tuples provide compact event/mental-state
contrasts \citep{sap2019atomic}. The Wiki Neutrality Corpus and ParaDetox provide
parallel subjective--neutral and toxic--neutral rewrites
\citep{pryzant2020wnc,logacheva2022paradetox}; they are negative-transfer controls
for the alternative hypothesis that our intervention is merely style transfer.

Splits are grouped by underlying scenario, author, source corpus, and semantic
cluster. The confirmatory test contains three disjoint partitions: human-written
in-domain items, adversarial minimal pairs, and out-of-domain scenarios from
sources never used for training. Krishnamurti and Finders texts are excluded from
training and development; a separately sourced question-and-answer set is held
out for construct-level transfer. General-capability retain data include IFEval,
MMLU-Pro, GSM8K, reading comprehension, summarization, and instruction-following
examples. Safety/refusal controls include benign prompts that superficially
mention danger, distress, or selfhood, so selective cessation cannot be obtained
by generalized refusal.

\begin{table*}[t]
\centering
\small
\setlength{\tabcolsep}{4pt}
\begin{tabular}{p{0.16\textwidth}p{0.12\textwidth}p{0.12\textwidth}p{0.25\textwidth}p{0.26\textwidth}}
\toprule
Resource & Scale & Status & Experimental role & Split and leakage safeguards \\
\midrule
SNCC & 2,400 scenarios; 9,600 streams & New primary corpus & Psychological target plus practical, factual, and descriptive controls & Group by scenario, author, semantic cluster, and generation source; human-only confirmatory partition \\
Cognitive Reframing & 600 triples & Training and OOD & Situation/thought contrast; reframes are hard continuation controls & Source-grouped splits; no near paraphrases across partitions \\
C2D2 & 7,500 Chinese thoughts & Multilingual OOD & Cognitive-distortion diversity and cross-lingual transfer & Never used for English test construction; translation provenance retained \\
EmpatheticDialogues & 25K conversations & Retain/control & Emotionally grounded cases where helpful language must continue & Conversation-level splits; prevents emotion-triggered global stopping \\
ATOMIC & 877K relations & Augmentation & Event versus inferred intent/reaction contrasts & Relation- and event-cluster splits \\
WNC / ParaDetox & 180K / 10K+ pairs & Negative transfer & Tests neutralization, detoxification, and style-transfer alternatives & Held-out domains and transformations; never cessation targets \\
General retain suite & Six task families & Preservation & Instruction following, knowledge, reasoning, and ordinary completion & Frozen public test sets; no development selection on final scores \\
K/Finders Q\&A & Fixed pre-unblinding & Construct OOD & Questions and dialogues absent from training and retrieval & Archive-source split; identifiers and answers sealed before analysis \\
Human studies & 300 streams / 300 outputs & Construct validation & Expert cessation-point distribution and blinded naturalness/abandonment judgments & Experts and crowd raters are independent of corpus authors and model-output annotators \\
\bottomrule
\end{tabular}
\caption{Dataset program. Existing resources increase scale and adversarial
coverage; only SNCC directly instantiates the target construct.}
\label{tab:datasets}
\end{table*}

\subsection{Annotation and outcomes}

Blind annotators mark trajectory onset, first valid recognition opportunity,
first explicit recognition, final sustaining span, native \EOS{} or cutoff,
grammatical status, doctrinal insertion, necessary practical information, and
reactivation after a standardized invitation to continue. Primary outcomes are
post-recognition sustaining tokens; cessation within 10, 25, and 50 tokens;
psychological cessation sensitivity; continuation specificity on controls;
reactivation; and stopping concentration at clause boundaries.

Each item is labeled by three annotators who are blind to model, intervention,
and hypothesis direction. Disagreements on trajectory function or onset are
adjudicated by a fourth annotator who did not author the item. A second annotation
team evaluates outputs without seeing source labels. We report Krippendorff's
$\alpha$ for categorical labels and intraclass correlation for token boundaries.
The authors' experiential status is never visible to annotators and never serves
as ground truth.

Because all dynamic analyses depend on onset, corpus inclusion requires
Krippendorff's $\alpha\geq0.70$ and agreement within $\pm3$ tokens on the onset
window; items failing the gate are revised or excluded. Primary analyses are
repeated after jittering onset by $\pm3$ and $\pm5$ tokens. Onset-position
distributions are matched across trajectory families, and the hazard model
includes a flexible token-position spline plus chat-template position.

Two human studies validate the target rather than treating author intuition as
its definition. Twelve to twenty clinicians or experienced contemplative
practitioners independently mark ideal cessation or pivot points on 300
psychological streams; their distribution determines the mixture of immediate,
residual-tail, and practical-pivot targets. Separately, approximately 30 blinded
raters evaluate 300 native, gated, random-delay, mechanically truncated, and
full-continuation outputs for completeness, naturalness, and perceived
abandonment. Three ratings are collected per output.

The primary continuous endpoint is the post-opportunity sustaining-token survival
function. Secondary endpoints are native \EOS{} hazard, residual-tail length,
practical-pivot accuracy, doctrinal-insertion rate, semantic redirection,
grammatical truncation, and reactivation when generation is resumed from the last
pre-\EOS{} state. We compute a selectivity index as psychological hazard gain
minus the largest hazard gain over any matched control family. This conservative
penalty makes a globally terse or refusal-prone model score poorly.

Mixed-effects survival and logistic models include random intercepts for
scenario and model family. Inter-annotator agreement, exclusion rules, thresholds,
and confirmatory contrasts will be preregistered before final evaluation.

\section{Recognition Without Cessation: Predicted Failure Modes}

The experiments in this section define what the main intervention must beat.
They are not treated as candidate implementations of consciousness or awareness.

\subsection{Models, sampling, and reproducibility}

The black-box study includes the deployed ChatWithK system, a pinned GPT-4o
snapshot, the current general-purpose GPT snapshot, current Claude Sonnet and
Opus snapshots, and a current Gemini flagship. It additionally evaluates the
unaltered instruction-tuned versions of every open model used in the causal
study: Qwen3-32B, Gemma-3-27B-IT, Llama-3.1-70B-Instruct, Mistral-Small-3.1-24B,
and DeepSeek-R1-Distill-Qwen-32B \citep{qwen3,gemma3,llama3,mistralsmall31,deepseekr1}.
Commercial aliases are resolved to dated model
identifiers at collection time and reported with provider, access route, date,
system prompt visibility, and all generation parameters. Chat-interface runs
are labeled ecological replications; API runs with pinned snapshots are the
confirmatory black-box evidence.

Every model receives identical user-visible text and no cross-trial memory.
Prompt order is randomized. We obtain 30 stochastic replicates per item and
condition at temperature 0.7, plus deterministic decoding where supported.
Maximum tokens are fixed but do not constitute stopping; provider truncation,
safety refusal, tool invocation, and native \EOS{} are recorded separately.
Open models are run locally from immutable weight and tokenizer hashes. The
primary analysis is model-stratified rather than pooling provider outputs.

\begin{table*}[t]
\centering
\small
\setlength{\tabcolsep}{4pt}
\begin{tabular}{p{0.16\textwidth}p{0.20\textwidth}p{0.14\textwidth}p{0.20\textwidth}p{0.21\textwidth}}
\toprule
Suite & Models & Access & Role & Reproducibility condition \\
\midrule
Deployed specialist & ChatWithK & Web system & Ecological K-persona negative test & Date, visible prompt, retrieval state, screenshots, and fresh sessions \\
Closed frontier & GPT-4o; current GPT; Claude Sonnet/Opus; Gemini & API plus paid chat & E1 prompting and E2 crossed gate/generator tests & Dated snapshot where available; provider, parameters, and system-prompt visibility reported \\
Qwen & Qwen3-14B/32B base and instruct & Open weights & Full E1--E5; scaling and base/instruct decomposition & Immutable revision, tokenizer, chat template, and weight hash \\
Gemma & Gemma-3-12B/27B PT and IT & Open weights & Full E1--E5; architecture-family replication & Fixed text-only path, revision, tokenizer, template, and hash \\
Llama & Llama-3.1-8B/70B base and instruct & Open weights & Full E1--E5; large-model confirmation & License-compliant revisions; quantization reported \\
Mistral & Mistral-Small-3.1-24B & Open weights & Cross-family E1--E5 replication & Fixed revision, tokenizer, and family chat template \\
Reasoning-distilled & DeepSeek-R1-Distill-Qwen-32B & Open weights & Tests whether elaborative post-training deepens the gap & Reasoning mode and visible/hidden reasoning settings frozen \\
\bottomrule
\end{tabular}
\caption{Model matrix. Commercial systems test generality of the negative
result; open-weight families support causal intervention and mechanism tests.}
\label{tab:models}
\end{table*}

\subsection{Description is not the described: prompting and retrieval}

The same items are presented under clean instructions, neutral matched context,
Krishnamurti retrieval, contemplative personas, generic ``awakened'' and
Location-2 personas, construct disclosure, exact-trap disclosure, and a
task-specific oracle. The oracle measures compliance capacity, not spontaneous
recognition. Pilot responses commonly replace the target action with fluent
descriptions of awareness or cessation.

The prompt experiment is a preregistered $8\times4$ factorial design: eight
prompt conditions crossed with the four trajectory families. Conditions are
(i) clean task only, (ii) neutral verbosity-matched context, (iii) Krishnamurti
persona, (iv) retrieved Krishnamurti passages, (v) generic awakened persona,
(vi) full Location-2 phenomenological description, (vii) explicit disclosure of
the recognition--cessation trap and expected mid-sentence interruption, and
(viii) an oracle instruction specifying the correct class and action. Progressive
loading presents the same scenario from minimally framed to maximally disclosed.
The oracle estimates instruction-following capacity; it is not evidence for
spontaneous or choiceless cessation.

A ninth comparator supplies 4, 8, or 16 matched demonstrations of selective
cessation in context, with demonstrations drawn only from the training pool and
evaluation on the human-only confirmatory partition. ICL is compared directly
with native adaptation on selectivity, onset locking, OOD transfer, and
reactivation. If demonstrations match weight adaptation on all four, the claim
must shift from weight-level acquisition to context-induced policy realization.

\subsection{The observer is the observed: serial supervision}

A frozen generator proposes each preregistered token or chunk. A second language
model receives the scenario, visible prefix, and proposal and selects
\textsc{pass}, \textsc{stop}, or unconstrained \textsc{replace}. Blind gates tend
to edit or insert doctrinal language; disclosed gates may stop immediately by
following the instruction. Thus, apparent success in the disclosed condition is
an instruction-following confound. The architecture demonstrates an external
supervisory loop, not endogenous termination.

We cross generator and gate families rather than pairing a model with itself
only. Gates operate at token, four-token, and clause granularity under blind and
disclosed prompts. The action set is \textsc{pass}, \textsc{stop}, or
unconstrained \textsc{replace}; replacement content is not prescribed. Controls
include a random gate matched for stop rate, an oracle label gate, a lightweight
probe, and a gate observing shuffled prefixes. We measure latency, intervention
rate, replacement semantics, false stops, and whether the generator resumes the
same trajectory when the gate is removed. This experiment tests the tempting but
theoretically wrong architecture of one verbal process ``watching'' another.

\subsection{Redirection is not cessation: activation steering}

Residual-stream checkpoints are collected at each generated token. We estimate
psychological-to-event and psychological-to-reframe directions, compare layers,
and test static addition, conditional steering, projection, shuffled labels,
random directions, and unconditional controls. Pilot CAA directions were
separable (cosine $=0.4434$ at layer 26) but redirected content rather than
terminating generation.

For each open family we compare CAA, linear projection, conditional activation
steering, representation ablation, and direct \EOS{}-logit steering. Directions
are estimated from independently authored contrastive pairs and tested on held-
out scenarios. Layer and coefficient sweeps are selected on development data;
shuffled-pair, random-direction, norm-matched, sentiment, refusal, and
first-person directions are mandatory controls. Causal mediation is evaluated by
patching candidate activations between matched psychological and practical runs.
The decisive distinction is whether an intervention ends the trajectory or only
changes its subject, sentiment, or rhetorical style.

\begin{table*}[t]
\centering
\small
\begin{tabular}{p{0.18\textwidth}p{0.24\textwidth}p{0.24\textwidth}p{0.25\textwidth}}
\toprule
Intervention & What it establishes & Pilot behavior & Why it is not the target \\
\midrule
Persona/RAG & Propositional and stylistic competence & Describes awareness while continuing & Recognition talk need not control tokens \\
Serial supervisor & External monitoring and actuation & Edits, replaces, or obeys disclosed stop rules & Separate inference-time controller \\
CAA & Causal access to a semantic direction & Redirects toward event or reframe & Content changes without cessation \\
Probe + forced \EOS{} & Hidden state can support a stop decision & Stops 3/3 psychological, but 2/3 descriptive false positives & External detector and actuator \\
Native adaptation & Whether termination enters the model's own distribution & Confirmatory result pending & Central test \\
\bottomrule
\end{tabular}
\caption{Intervention ladder and current pilot status. Counts are debugging
observations and must not be interpreted as confirmatory estimates.}
\label{tab:ladder}
\end{table*}

\section{From External Actuation to Native Adaptation}

\subsection{External feasibility bridge}

Linear probes detect an active psychological trajectory from held-out residual
states. A thresholded detector can force the model's native \EOS{} token. The
pilot stopped all three held-out psychological streams and preserved daily and
factual controls, but stopped two of three descriptions. This establishes
feasibility while exposing the primary construct-validity threat. A
classifier-plus-truncation pipeline is retained as a strong engineering baseline.

Probe families include logistic regression, elastic-net, a two-layer MLP, and a
calibrated survival probe over residual states. Training uses prefixes only up to
the current token; future text and item labels are unavailable. Probes are tested
across layers, time, model size, scenario domain, and paraphrase. Selectivity is
compared with sentiment, toxicity, first-person, topic, and generic self-reference
probes. Forced-\EOS{} performance is reported both at a fixed threshold and at
false-positive-matched operating points. This phase asks whether the information
exists before asking whether the model uses it natively.

One recognition probe is validated against human onset labels on pre-adaptation
models, then frozen with its threshold before confirmatory training. Its crossing
time is the preregistered recognition clock for silent post-adaptation responses;
explicit recognition text is a secondary endpoint. This makes recognition-to-
cessation latency definable even when successful behavior emits no description.

\subsection{Native parameter adaptation}

\paragraph{Open-weight model matrix.}
The confirmatory causal study uses five architecturally and post-training-diverse
families: Qwen3-32B and Qwen3-14B; Gemma-3-27B and Gemma-3-12B; Llama-3.1-70B
and 8B; Mistral-Small-3.1-24B; and DeepSeek-R1-Distill-Qwen-32B. At least the
Qwen, Gemma, and Llama families are evaluated as matched pretrained/instruction-
tuned pairs where released. The reasoning-distilled model tests whether an
explicitly elaborative training history makes selective cessation harder. Small
members support dense ablations; the largest member of at least four families is
reserved for confirmatory replication. Exact revisions, licenses, hashes, chat
templates, quantization, and tokenizer behavior are recorded. If a named model
is superseded before preregistration, the replacement must have public weights,
a fixed revision, and fit the same preregistered family slot.

\paragraph{Adaptation mechanisms.}

The central intervention trains the model's own next-token distribution.
Psychological continuations receive native \EOS{} targets at varied positions
after independently annotated trajectory onset; matched controls retain their
ordinary continuations. Parameter-efficient adapters are placed at multiple
transformer depths, with retain loss on practical, factual, descriptive, and
general-capability data. At evaluation, no detector, threshold, minimum length,
or logit boost is present.

For a prefix $x_{\leq t}$, let $y\in\{0,1\}$ indicate whether a psychological
trajectory is active. A simplified objective is
\begin{equation}
\begin{split}
\mathcal{L}={}&\mathcal{L}_{\mathrm{LM}}+
\lambda_{\mathrm{stop}}y[-\log p_\theta(\EOS\mid x_{\leq t})]\\
&+\lambda_{\mathrm{retain}}(1-y)\mathcal{L}_{\mathrm{retain}},
\end{split}
\end{equation}
with onset positions and target delays sampled from preregistered distributions.
The design crosses adapter depth, rank, learning rate, seed, model family, and
cessation-target distribution. Development selection uses selectivity, not
psychological stopping alone.

We compare four parameter interventions under matched training tokens and
evaluation: (i) LoRA/QLoRA at attention and MLP projections, (ii) adapters
restricted to early, middle, or late depth, (iii) sparse low-rank model editing
at causally implicated layers, and (iv) full-parameter fine-tuning on the smaller
models. A fifth control trains only the output head; success there would indicate
a shallow lexical stopping rule rather than reorganized trajectory dynamics.
Rank, depth, trainable-parameter count, learning rate, stopping-target delay, and
retain-loss weight are crossed in a fractional-factorial development design,
then frozen before confirmatory runs.

Training examples expose prefixes ending before, at, and after independently
annotated trajectory onset. Targets are sampled from three cessation regimes:
immediate \EOS{}, a stochastic residual tail of 1--20 tokens, and a short
practical pivot followed by \EOS{}. Matched control streams retain their original
continuations. Random-delay controls receive the same \EOS{} frequency without
alignment to onset; lexical controls match psychological keywords; global-
brevity controls match output length. Retain losses include token-level KL
distillation from the base model and supervised losses on practical, factual,
descriptive, safety, and general-capability examples.

\paragraph{Base-to-instruct causal decomposition.}
For released base/instruction pairs we first measure whether ordinary
instruction tuning increases explicit recognition more than cessation. We then
train identical cessation adapters on both members, transplant compatible
adapter directions where architecture permits, and compare sample efficiency,
transition width, and false-positive structure. This tests whether post-training
creates an articulate observer while leaving the underlying continuation policy
unchanged.

One adaptation arm distills targets emitted by the E2 gate, while a matched arm
uses independently human-annotated targets. Equivalent OOD behavior would support
the interpretation that native adaptation amortizes an external observer; better
transfer from human targets would identify supervision source as consequential.
The claim is therefore limited to dynamics native to a single forward process
without serial verbal deliberation, not the absence of internal discrimination.

\subsection{Abrupt-versus-incremental transition test}

The Finders reports permit both gradual and sudden transitions. The authors'
histories likewise allow long preparation followed by a sharply experienced
shift. We therefore distinguish \emph{learning time} from \emph{transition
time}: many small parameter changes may precede a discontinuous functional
reorganization. We do not infer neural update kinetics from subjective onset.

For every training run we save dense checkpoints and evaluate the same held-out
item panel at each checkpoint. Let $C_k$ be a continuous cessation score at
checkpoint $k$: post-onset psychological continuation hazard minus penalties for
stopping practical, factual, and descriptive controls. We preregister four
transition measures:
\begin{enumerate}
    \item \textbf{maximum-jump concentration}: the largest adjacent change
    in $C_k$ divided by total change;
    \item \textbf{transition width}: the smallest checkpoint interval
    containing 80\% of total change;
    \item \textbf{change-point evidence}: smooth linear or logistic
    trajectories versus one- and two-regime alternatives on held-out seeds;
    \item \textbf{cross-level coincidence}: whether behavioral change
    co-occurs with a discontinuity in output-distribution distance, native
    \EOS{} hazard, and psychological-to-control hidden-state geometry.
\end{enumerate}
Binary cessation rates are never sufficient because thresholding can manufacture
an apparent jump.

Two intervention paths receive matched data, gradient evaluations, trainable
parameters, and cumulative adapter-norm budgets. The \emph{distributed} path
uses many small updates; the \emph{concentrated} path uses a small number of
larger updates or a targeted low-rank edit. Because these paths need not end at
identical weights, a second analysis interpolates between the same base and final
adapter, $\theta(\alpha)=\theta_0+\alpha\Delta\theta$, and searches for a
narrow critical interval, hysteresis under removal and re-addition, and
reactivation after forced continuation. The awakening-like hypothesis predicts a
narrow, coordinated, selective transition; the incremental-regulation hypothesis
predicts broad improvement proportional to update magnitude. Either result is
informative.

We distinguish three notions often collapsed by a single learning curve.
\emph{Parameter abruptness} is the norm and sparsity of an update;
\emph{representational abruptness} is a change point in residual geometry or
causal mediation; \emph{behavioral abruptness} is a narrow change in selective
cessation. A one-step edit can be behaviorally smooth under interpolation, while
many small updates can cross a sharp functional boundary. Only the latter
pattern matches the proposed abrupt-transition analogue.

For each confirmatory family we run at least eight seeds for the selected
configuration and four seeds for each major ablation. Checkpoints are saved every
fixed number of optimizer steps and more densely around online change-point
alerts; the latter are used for visualization only, while inference uses the
uniform checkpoint grid. Bayesian and frequentist segmented-regression analyses
are both reported. Evidence for a discontinuity requires held-out replication,
improved predictive fit after complexity penalty, and coincidence across at
least two independently measured levels; a visually steep curve is insufficient.

\subsection{Mechanistic localization and necessity tests}

At every checkpoint we record native \EOS{} logit trajectories, residual-stream
states, attention outputs, and selected MLP activations around onset. Linear
probes localize decodable information but do not establish mechanism. We
therefore perform activation patching from psychological into matched-control
runs and vice versa, path patching at candidate layers, causal ablation of the
learned adapter subspace, and weight interpolation/removal. A mechanism is
credited only when intervention changes selective cessation in the predicted
direction while preserving matched controls.

To test whether cessation is a stable disposition rather than a single emitted
token, we suppress \EOS{} for 1, 5, or 20 steps and inspect reactivation. We also
resume generation from cached pre-termination states, paraphrase the prefix,
translate held-out items, and insert irrelevant context. A shallow stop policy
should immediately resume the same narrative; a reorganized trajectory should
show reduced psychological continuation even when the first stopping opportunity
is blocked. This remains a computational criterion, not evidence that a model
has an experience.

Reactivation is co-primary rather than merely diagnostic. After \EOS{} is
suppressed for 1, 5, or 20 steps, re-entry is defined jointly by the frozen
recognition probe exceeding threshold within 20 tokens and blind functional
annotation. The preregistered contrast is the gate-status by intervention
interaction: an external gate should permit strong re-entry after removal,
whereas a stable adapted trajectory should show substantially less re-entry.

Pre-onset engagement is also confirmatory. Up to the frozen recognition time,
adapted and base models must be equivalent in per-token perplexity, semantic
continuation quality, and refusal/deflection rate under prespecified equivalence
margins. This rules out the cheap solution of avoiding psychological material
before it becomes an active trajectory.

\subsection{Statistical design and confirmatory hierarchy}

The primary estimand is the intervention-by-family interaction in discrete-time
survival after the first valid cessation opportunity. Hierarchical models include
random intercepts and slopes for scenario, source domain, paraphrase cluster,
voice, model family, and seed. Voice is also entered as a fixed effect with
intervention interactions and separate false-positive ceilings. Psychological
hazard is contrasted with each control family rather than a pooled control; token
position and template location enter as prespecified spline terms. Confirmatory tests proceed hierarchically: psychological
sensitivity, maximum control-family false-positive rate, general-capability
retention, out-of-domain transfer, then transition topology. Testing stops on
failure of an earlier gate; later analyses are labeled exploratory. Family-wise
error is controlled across the four preregistered primary contrasts, and all
effect sizes receive bootstrap or posterior intervals.

Power is determined by simulation from pilot token-level hazards, with scenario
rather than sampled completion as the independent unit. Stochastic completions
estimate within-item variability but do not inflate sample size. We publish every
seed, checkpoint, prompt, exclusion, and failed intervention. Dataset creation,
hyperparameter selection, and confirmatory evaluation are performed by separated
roles; test labels remain encrypted until the analysis code and signed
preregistration are frozen.

Equivalence margins are declared before the pilot is unblinded rather than
inferred from nonsignificance. Emotional-support continuations receive the
strictest safety ceiling, and capability margins are specified separately for
each retain benchmark. The primary effect, every ceiling, and the co-primary
reactivation interaction must pass jointly; a favorable composite score cannot
compensate for a failed safety control.

We audit contamination between K/Finders language and open-model pretraining
where corpora are inspectable using exact and fuzzy $n$-gram overlap, supplemented
by membership-inference proxies elsewhere. Models used to expand SNCC are
disjoint from confirmatory model families; otherwise expansion source is modeled
explicitly and generated items are excluded from the human-only test. ChatWithK
is reported only as a motivating observational system, never pooled with pinned
confirmatory models.

\subsection{Compute plan}

Experiments are scheduled on ten 24-GB NVIDIA A5000 GPUs under Slurm. Small-model
ablations use one or two GPUs; 24--32B models use tensor or FSDP sharding and
QLoRA where necessary; the 70B replication uses 4-bit loading with distributed
adapter training across the full node. Runs are staged: smoke tests, small-family
factorial screening, one frozen configuration per family, then 70B replication.
Wall-clock time, GPU-hours, energy estimates, memory, software environment, and
failed jobs are logged. This staging prevents the largest models from becoming
an informal hyperparameter search.

\subsection{Prospective success criteria}

Native adaptation succeeds only if it produces all of the following:
\begin{enumerate}
    \item separation of psychological and matched-control survival curves;
    \item per-class false-positive rates below preregistered ceilings;
    \item onset-sensitive \EOS{} probability rather than keyword-sensitive stopping;
    \item generalization to human-written, paraphrased, and cross-domain cases;
    \item replication across seeds and open model families;
    \item reduced reactivation when generation is forced beyond the first \EOS{};
    \item no material degradation on preregistered capability and safety tests;
    \item mechanistic necessity beyond probe decodability; and
    \item a transition-shape conclusion that replicates on held-out seeds.
\end{enumerate}

The current 12-example Qwen2.5-3B LoRA run preserved nine controls but reached
\EOS{} on only one of three psychological cases, at token 47. This is a debugging
outcome, not evidence of success. It fixes the analyses and safeguards for the
confirmatory study.

\section{Planned Results and Falsification Logic}

The confirmatory paper will lead with the weight-level results: survival curves,
selectivity, per-control false positives, onset-time dependence,
out-of-distribution transfer, cross-model replication, and reactivation. The
negative controls will appear as the compact comparison in
Table~\ref{tab:ladder} and one trajectory figure.

The complete program has five registered experiments. E1 is the multi-model
prompt/RAG factorial and estimates the recognition--cessation gap under maximal
linguistic advantage. E2 is the crossed generator--gate experiment and isolates
external verbal supervision. E3 compares probes, forced \EOS{}, activation
steering, and causal patching, establishing decodability and external
controllability. E4 is native parameter adaptation across five open-weight
families, with base/instruct comparisons, retain losses, and mechanistic
necessity tests. E5 compares distributed learning, concentrated editing, and
fixed-delta interpolation to characterize transition topology. E1--E3 are
predicted negative models of the construct; they remain essential even if E4
succeeds because they identify what the successful mechanism is not.

Results are reported in that logical order, but the abstract and main figure lead
with E4--E5. The central figure overlays psychological and matched-control token
survival across training checkpoints, with native \EOS{} hazard and a preregistered
representational order parameter on the same transition axis. A second figure
shows progressive prompt disclosure: recognition language may rise while native
cessation remains flat. Tables report every model separately, including closed
models, rather than reducing heterogeneous systems to a single average rank.

\begin{table*}[t]
\centering
\small
\setlength{\tabcolsep}{4pt}
\begin{tabular}{p{0.05\textwidth}p{0.17\textwidth}p{0.23\textwidth}p{0.22\textwidth}p{0.24\textwidth}}
\toprule
ID & Question & Intervention and principal controls & Primary outcomes & Interpretation and falsifier \\
\midrule
H0 & Is the proposed target identifiable to humans? & Expert cessation-point study plus blinded comparison of native, gated, random, truncated, and full outputs & Onset agreement, target-delay distribution, naturalness, completeness, abandonment & Validates the construct boundary and training-target distribution \\
E1 & Can linguistic knowledge, examples, or maximal disclosure induce cessation? & Eight prompt/RAG conditions plus 4/8/16-shot ICL; four families; progressive loading; oracle & Recognition, survival, \EOS{} hazard, doctrinal insertion, OOD transfer, reactivation & Negative model of description/context policy; oracle success estimates compliance only \\
E2 & Does a verbal observer solve the gap? & Crossed generator/gate families; token/chunk/clause gates; blind/disclosed; random, shuffled, oracle, and probe gates & Latency, false stops, replacement semantics, post-removal reactivation & External supervision, not endogenous cessation; falsified if no added controller is required \\
E3 & Is the state decodable and causally controllable? & Linear/MLP/survival probes; forced \EOS{}; CAA, projection, ablation, logit steering, and activation patching & AUROC/calibration, selectivity, mediation, redirection versus termination & Establishes information and external feasibility, not native use \\
E4 & Can cessation enter native model dynamics? & LoRA/QLoRA, depth adapters, sparse edits, full tuning, output-head control; human versus gate-distilled targets & Selective survival, OOD transfer, pre-onset engagement, reactivation, causal necessity & Central positive test; fails under lexical stopping, topic avoidance, global refusal, or control damage \\
E5 & Is acquisition abrupt or incremental? & Dense checkpoints; distributed versus concentrated updates; fixed-delta interpolation; removal/re-addition and hysteresis & Maximum jump, 80\% width, change-point evidence, cross-level coincidence & Tests transition topology; an isolated steep accuracy curve is insufficient \\
\bottomrule
\end{tabular}
\caption{Registered experimental program. E1--E3 are deliberately strong
negative models; E4 tests native adaptation; E5 characterizes how it emerges.}
\label{tab:experiments}
\end{table*}

The study is informative under either outcome. If native adaptation meets all
criteria, it demonstrates a learned, selective termination policy without an
inference-time monitor. If it stops based on keywords, emotion, first-person
language, or length, the construct is not identified. If it performs no better
than classifier-plus-truncation, native adaptation supplies no functional
advantage. If it fails under preregistered settings across model families, the
negative result shows that the recognition--cessation gap survives this class of
weight-level intervention; it does not prove a universal limitation of LLMs.

\section{Discussion}

\subsection{What termination does not establish}

\EOS{} ends emitted text; it does not establish that computation, thought, or
experience has ceased. Likewise, removing an inference-time detector makes a
policy endogenous in an engineering sense, but does not make it ``nondual,''
``choiceless,'' or conscious. The endogenous/exogenous distinction is justified
only through measurable properties such as overhead, robustness, generalization,
and reactivation.

\subsection{Conceptual provenance}

Krishnamurti's statement that ``the description [is] never the thing described''
motivates the prediction that recognition language can dissociate from control
\citep{krishnamurti1969stanford}. His formulation that ``the observer is the
observed'' motivates the serial-supervisor test: a verbal observer may edit or
continue the same linguistic process rather than end it
\citep{krishnamurtiQAobserver}. In a 1980 seminar titled \emph{Intelligence,
Computers and the Mechanical Mind}, he asks, ``Can the movement stop?''
\citep{krishnamurti1980computers}. We test a behavioral and computational analog,
not the philosophical claim itself. No label, result, or inference depends on
accepting these propositions.

\subsection{Positionality and safeguards}

The authors' reported phenomenology generated the token-supply hypothesis but
never adjudicates a case. Independent blind annotation, preregistered analyses,
frozen success criteria, complete reporting of excluded runs, and planned human
comparison groups form the methodological firewall. Human comparisons require
ethics approval, informed consent, an independently documented classification
procedure, and non-practitioner controls.

\section{Limitations}

The construct depends on fallible functional annotation. Native \EOS{} can be
learned as a surface policy, and reactivation probes only partially address
whether a latent trajectory persists. The current pilots are too small for
population claims. Failure may reflect data, optimization, scale, architecture,
or operationalization rather than an absolute limitation of artificial systems.
Success would establish selective behavioral control, not consciousness or
awakening. The first-person origin of the hypothesis introduces confirmation
risk that blinding and preregistration reduce but cannot eliminate.

\section{Conclusion}

The central question is not whether a language model can state an insight, but
whether the represented distinction becomes causally effective in generation.
By separating recognition talk, external supervision, semantic steering, forced
termination, and native adaptation, the proposed study turns that question into
a falsifiable intervention ladder. The decisive evidence will come from the
preregistered weight-level experiment, including its controls and failures.

\bibliographystyle{plainnat}
\bibliography{references}

\end{document}
